% Working note, not thesis prose -- current state, to revisit together once confirmed.
% Numbers and figure from scripts/survey_topology.py and scripts/make_tree_figure.py, 2026-09-14.

\section*{The coronary tree graph (current approach)}

Before I can measure anything about a coronary tree's shape, I need the tree as an actual data
structure, not just a pile of points -- something with a notion of "these two branches meet here,"
"this is upstream of that." Every topology feature extraction method and analysis relies on the
following:

I start by building the tree graph from the centerline files ImageCAS-X ships, one per side of the
heart -- these are the dataset's ground truth, not a prediction from any model. No segmentation
network is involved anywhere in this section; that only shows up later, in the Herlev--{\O}sterbro
part.

That said, "ground truth centerline" needs unpacking, because it isn't a human-drawn line. The
\emph{mask} is ground truth in the usual sense: analysts corrected it by hand, slice by slice. But
the centerline sitting on top of that mask is not hand-drawn -- the original hand-traces were
thrown away and replaced by running a fixed, deterministic procedure over the corrected mask:
skeletonize it to a 1-voxel-wide line, then Gaussian-smooth it ($\sigma = 0.5$\,mm)
\cite{bransby2026imagecasx}. So even in the ground-truth data, the centerline's geometry and
connectivity come from that fixed algorithm, not from a person tracing the vessel; only the segment
names (LM, LAD, \ldots) are human, carried over from the old hand-traces by nearest-distance
matching. Worth keeping straight: "algorithmic" here means "produced by that fixed skeletonization
step," not "predicted by a network" -- the mask underneath it is still the corrected ground truth.

As a file, a centerline is just points plus a bunch of polylines connecting them, with no notion of
direction and no notion of "this is one vessel." That's not enough on its own -- you can't get a
branch's direction out of it without first sorting out which way is outward, and a single vessel is
often chopped across several polylines. So \texttt{topology/graph.py} does the actual building:
root each side at its ostium, orient every polyline outward from there, glue the pieces of one
vessel back into a single segment, and name each segment by majority vote over its interior points
(the junction point itself gets skipped, since it still carries the parent vessel's label).

I checked this against the flags the dataset ships (which points are ends, which are bifurcations):
across all 1600 sides, the graph's connectivity matches exactly, and 1584 sides come out as one
clean tree. Not fully independent evidence -- the flags were defined by the same degree rule -- but
it confirms the connectivity was recovered right.

Here's what it actually looks like for one case (\Cref{fig:week3-tree-case1}): on the left, the
centerline as it sits in space, with the five bifurcation angles marked; on the right, the same
tree redrawn as a dendrogram, height being distance from the ostium, so branch order and segment
length are easy to read off directly instead of squinting at 3D.

% Copy figures/1_05_tree.png (this repo) to content/background/figures/coronary_tree_case1.png
% when pasting this into the real thesis, same as the other week files' figures.
\begin{figure}[htbp]
  \centering
  \includegraphics[width=\linewidth]{content/background/figures/coronary_tree_case1.png}
  \caption{Case 1's left and right coronary trees: the centerline in space (left) and the same tree
    as a dendrogram (right two panels), colored by artery. Stars mark the five bifurcations this
    week's feature measures.}
  \label{fig:week3-tree-case1}
\end{figure}

\textbf{Still open: getting this onto Herlev--{\O}sterbro.} None of this exists there yet -- no
mask, no centerline, no ostium, no names. Rough plan: predicted mask from the segmentation network,
skeletonize it the same way, ostium from a TotalSegmentator aorta
\cite{wasserthal2023totalsegmentator}, and names from... something, probably training on the 14
labels directly instead of vessel-vs-background. That last part I don't have a good answer for yet.
Before touching a real Herlev--{\O}sterbro scan, the plan is to run the whole thing on the 160
ImageCAS-X test cases first and see how far the angles drift once the tree comes from a predicted
mask instead of the corrected one -- that drift is the error this extraction adds, and it needs to
be small next to how much the angle varies between people.
